\documentclass[a4paper,14pt]{extarticle}
\usepackage{styles}

\begin{document}

\litSubsubsection{2.2.2 GPU Occupancy and Utilization}

Device occupancy is a key factor in GPU effectiveness. This section reviews it as a metric for predicting when parallel GAs will outperform sequential implementations.

\textbf{Theoretical Occupancy.} ``Occupancy is defined as the ratio of active warps on an SM to the maximum number of active warps supported by the SM.''~\cite{nvidia2024occupancy}:

\begin{equation}
\text{Occupancy} = \frac{\text{Active Warps per SM}}{\text{Maximum Warps per SM}}
\label{eq:occupancy}
\end{equation}

Modern GPUs (CC 8.x) support 48--64 warps per SM (1536--2048 threads)~\cite{jia2018volta}. Higher occupancy enables better latency hiding~\cite{volkov2010occupancy}. Table~\ref{tab:threads_per_sm} shows values by compute capability.

\begin{table}[H]
\centering
\small
\begin{tabular}{@{}lccc@{}}
\toprule
\textbf{CC} & \textbf{Example GPUs} & \textbf{Threads/SM} & \textbf{Warps/SM} \\
\midrule
7.5 & RTX 2080, T4, GTX 1660 & 1024 & 32 \\
8.0 & A100 & 2048 & 64 \\
8.6 & RTX 3050, RTX 2050 & 1536 & 48 \\
8.9 & RTX 4090, RTX 4080 & 1536 & 48 \\
\bottomrule
\end{tabular}
\caption{Maximum threads per SM by compute capability~\cite{nvidia2024programming}.}
\label{tab:threads_per_sm}
\end{table}

SM count can be queried at runtime via \texttt{cp.cuda.runtime.getDeviceProperties()} in CuPy~\cite{cupy_docs}. Achieved occupancy may be lower than theoretical due to workload imbalance or insufficient parallelism~\cite{nvidia2024occupancy}.

\subsubsection*{Population Size and GPU Utilization}

Each GA individual maps to one thread. For full utilization, the population must fill all SMs:

\begin{equation}
P_{\min} = \text{SM Count} \times \text{Threads per SM} \times \text{Target Occupancy}
\label{eq:pmin}
\end{equation}

When population size is too small, SMs remain idle. ``The smaller the benchmark and population are, the slower a GPU will be compared to a multicore CPU''~\cite{jaros2012fair}, since the GPU relies on thousands of concurrent threads to hide memory latency~\cite{pospichal2010parallel,carpentries_gpu}.

Occupancy is further limited by per-thread resource consumption: complex fitness functions consume more registers (reducing concurrent warps), and caching data in shared memory reduces room for concurrent blocks. However, maximum occupancy does not always yield maximum performance - for memory-bound kernels, performance plateaus at moderate occupancy levels~\cite{volkov2010occupancy}.

\textbf{Conclusions for Genetic Algorithm Implementation.} Population sizes should exceed a minimum threshold to justify GPU overhead. Block sizes of 128--256 threads generally balance occupancy and register pressure~\cite{nvidia_occupancy}. Populations below $\sim$1000 individuals are unlikely to benefit from GPU acceleration on modern hardware. These constraints inform the breakeven condition in Section 2.2.3.

\end{document}
